% <<<PDFPAGE 46>>>
\chapter{可观测性与生产工程}

\section{生产推理技术栈}

大语言模型通过离线基准测试，只是部署到生产环境的一步。生产环境要求系统具备可观测性（observability）和可靠性，还需要管理成本并持续改进。生产推理技术栈在服务引擎之上增加了以下几层。
\begin{itemize}
  \item \textbf{日志与追踪（logging and tracing）}。记录每个请求的 TTFT、TPOT、词元数、模型版本和错误状态。
  \item \textbf{指标与告警（metrics and alerting）}。通过仪表板跟踪 p50、p95、p99 时延，以及吞吐量、GPU 利用率和错误率。
  \item \textbf{A/B 测试基础设施（A/B testing infrastructure）}。在不同模型版本之间分配流量，让新版本逐步接收生产请求。
  \item \textbf{容量规划（capacity planning）}。预测查询量，并据此提前配置硬件。
  \item \textbf{成本归属（cost attribution）}。按团队或功能拆分词元成本。
\end{itemize}

\section{关键指标详解}

\textbf{首词元时延（TTFT）}是从收到请求到向用户返回第一个生成词元的时间，主要取决于预填充计算。用户感受到的 TTFT 是模型的“思考时间”。如果 TTFT 超过 1~s，即使后续词元输出很快，用户也会觉得应用响应迟缓。

\textbf{每输出词元耗时（TPOT）}是流式输出阶段相邻两个输出词元之间的平均时间间隔。人每分钟大约能读 250 个英文单词（约合每秒 5 个词元），所以 TPOT 小于 200~ms 时，流式输出通常比较流畅。TPOT 主要取决于解码步的时延和 GPU 的访存带宽。

\textbf{有效吞吐量}是同时达到全部 SLO 的请求所占比例，这里要求 TTFT 与 TPOT 两项均达标。如果 TTFT 和 TPOT 各有 95\% 的请求达标，并且两项达标事件相互独立，那么有效吞吐量为 $0.95^2\approx90\%$。相比单项指标的达标率，有效吞吐量提供了更多信息。

\textbf{模型浮点运算利用率（MFU）}是实际用于模型计算的 FLOP/s 占理论峰值 FLOP/s 的比例。经过充分优化的系统在预填充阶段能达到 30\%–60\% 的 MFU。解码阶段受访存带宽限制，MFU 要低得多。MFU 是理解硬件使用效率的一项重要指标。

\section{故障模式与排查}

% <<<PDFPAGE 47>>>
\begin{center}
\begin{tabularx}{.94\linewidth}{
  @{}>{\raggedright\arraybackslash}p{.23\linewidth}
  >{\raggedright\arraybackslash}X
  >{\raggedright\arraybackslash}X@{}
}
\toprule
\textbf{现象} & \textbf{可能原因} & \textbf{排查方法} \\
\midrule
TTFT 突然升高 & 排队时间过长、预填充抢占 & 监测队列深度、分解预填充时延 \\
TPOT 恶化 & KV 缓存淘汰、批大小增大 & 查看 KV 缓存命中率与解码步时延 \\
显存不足（Out of Memory，OOM） & KV 缓存溢出、请求激增 & 监控显存、查看 PagedAttention 块统计 \\
吞吐量骤降 & 张量并行中的通信瓶颈 & 对照 GPU 利用率与 NVLink 带宽 \\
输出质量下降 & 模型版本有问题、采样参数发生变化 & 查看日志中的模型版本标签、做在线评估 \\
\bottomrule
\end{tabularx}
\end{center}

\section{成本建模}

弄清每百万输出词元的成本并设法降低它，对业务能否持续至关重要。主要成本按下式计算。
\[
\text{每百万词元的成本}
=\frac{\text{每小时 GPU 成本}\times\text{使用小时数}}{\text{生成词元总数}}
\]
\translatornote{上式沿用原书。若分母是生成的词元总数，右端所得为每词元成本。计算每百万词元的成本时，需再乘以 $10^6$。}

可以从以下方面降低成本。
\begin{enumerate}
  \item \textbf{提高 GPU 利用率}。每批纳入更多请求，使用连续批处理。
  \item \textbf{减少模型计算量}。使用量化、蒸馏或更小的模型。
  \item \textbf{使用竞价实例或可抢占实例（spot/preemptible instances）}。此类实例比按需实例便宜 60\%–80\%，批处理作业通常可以接受。
  \item \textbf{选择合适的硬件}。对受访存带宽限制的工作负载，较早一代的 GPU（如 A100 相比 H100）可能更划算。
  \item \textbf{减少输出词元数}。通过指令微调（instruction tuning）让模型回答得更简洁。
\end{enumerate}

\begin{interview}
\question{1}{生产环境中的大语言模型服务出现 TTFT 的 p99 突然升高，你会如何排查？}
我会按层逐步缩小排查范围。（1）\textbf{网络层}。检查 p99 的上升是否与请求到达率上升相关。如果相关，系统可能在排队。再查看队列深度指标。（2）\textbf{预填充计算}。把预填充时延和排队等待时间分开。如果预填充时延突增，检查输入提示的长度分布是否改变，因为较长的提示需要更长的预填充时间。（3）\textbf{资源争用}。检查 GPU 利用率和 NVLink 带宽。如果同一硬件上还部署了其他工作负载，它们可能在争用 GPU 资源。（4）\textbf{KV 缓存压力}。如果 PagedAttention 分配块的速率上升，可能是请求吞吐量增加，造成计算资源争用。（5）\textbf{模型版本或配置变更}。对照日志，
% <<<PDFPAGE 48>>>
看同一时段是否有部署操作。（6）\textbf{硬件故障}。性能下降的 GPU 或 NVLink 链路可能使吞吐量降低。可用 \texttt{nvidia-smi} 查看错误计数。我会将这些情况与时延突增的时间点对应起来，向前追查原因。

\question{2}{什么是有效吞吐量？为什么它比原始吞吐量或单项指标的 SLO 达标率更有意义？}
吞吐量衡量每秒生成的词元总数，单凭它看不出用户的体验是否达标。系统可以把大量长请求放在同一批次里以获得很高的吞吐量，同时仍有请求的 TTFT 超出 SLO。单项 SLO 的达标率能多说明一些情况，例如“95\% 的请求达到 TTFT 的 SLO”。不过，各项 SLO 往往分别衡量，这会掩盖多项指标的共同达标情况。有效吞吐量是同时满足所有 SLO 的请求所占比例，包括 TTFT、TPOT 和其他已设定的目标。它反映系统真正达到用户体验要求的程度。如果 TTFT 和 TPOT 的达标率都是 90\%，并且未达标事件相互独立，有效吞吐量可能低至 81\%。做调度决策时，我会以有效吞吐量为优化目标。考虑 SLO 的调度器应尽量提高有效吞吐量，单纯提高总吞吐量并不够。

\question{3}{如何设计 A/B 测试框架，在生产环境中比较两个大语言模型版本？}
我会这样设计。（1）\textbf{流量划分}。对用户 ID 做一致性哈希（consistent hash），从而确定性地分配流量。同一用户始终访问同一个模型版本，可减少用户偏好差异造成的混杂。常见做法是初期按 $90{:}10$ 分流，随测试结果逐渐明朗而扩大测试流量。（2）\textbf{指标}。跟踪两个版本的时延（TTFT、TPOT）、质量信号（赞成或反对反馈、会话长度、后续提问）以及单请求成本。（3）\textbf{统计显著性}。采用序贯 A/B 测试，例如序贯概率比检验（Sequential Probability Ratio Test，SPRT），事先设定第一类和第二类错误的界限，据此判断何时已有足够的数据。（4）\textbf{防护规则}。配置自动回滚条件。如果测试版本的安全违规、错误率或时延恶化超过阈值，自动将全部流量切回对照版本。（5）\textbf{日志}。为每条请求和响应记录实验版本 ID、模型版本和各环节的完整耗时，以便事后分析。（6）\textbf{发布阶段}。从不影响用户的影子模式（shadow mode）开始，依次扩大至 1\%、10\%、50\% 和 100\%。

\question{4}{某大语言模型服务预计在六个月内迎来十倍的用户规模，你会如何做容量规划？}
（1）\textbf{测量现状}。记录当前每秒查询数（Queries Per Second，QPS）、平均输入和输出词元数，以及 p95 时延。（2）\textbf{预测 QPS}。按 QPS 增长到十倍估算，并考虑季节性和高峰、低谷时段。（3）\textbf{建立吞吐量模型}。根据当前 GPU 数量和单卡吞吐量，计算高峰时段所需的 GPU 小时数，并为突发流量预留 30\%–50\% 的余量。（4）\textbf{优化成本}。在增加服务器数量之前，尽量优化单机。可否用 4 位量化将单卡吞吐量翻倍？可否升级到访存带宽约为 A100 两倍的 H100？能否为系统提示启用前缀缓存？（5）\textbf{自动扩缩容（auto-scaling）}。用 Kubernetes 或云平台原生能力实现自动扩缩容，并为 GPU 资源准备留出时间。GPU 实例初始化需要 2–10 分钟。（6）\textbf{多区域部署}。提前规划地理分布，以降低时延并避免受限于单个区域的容量。做预算时，假定每张 H100 的成本为 2 美元/小时，充分利用时每小时可提供约 100 万个输出词元。

\question{5}{如何测量并提高模型浮点运算利用率（MFU）？}
% <<<PDFPAGE 49>>>
\textbf{测量 MFU}。先计算前向传播的浮点运算量。对于具有 $N$ 个参数的 Transformer，按每词元约 $6N$ 次运算估算，乘以每秒处理的词元数，再除以 GPU 的峰值 FLOP/s。使用 NVIDIA Nsight Compute 时，可以测量“SM Throughput”占峰值的比例。典型的 MFU 在预填充阶段为 30\%–45\%（受 GEMM 限制），在解码阶段低于 10\%（受访存带宽限制）。

\textbf{提高 MFU}。在预填充阶段，（1）用 FlashAttention-3 提高注意力内核的利用率。（2）分析性能并融合层归一化与残差运算。（3）大批次使用跨 GPU 的张量并行。（4）令批大小为张量核心分块尺寸的整数倍（如 16 的倍数）。在解码阶段，MFU 本来就低，更合适的指标是访存带宽利用率，目标可定在 HBM 峰值带宽的 80\% 以上。改善办法包括（1）增大批大小以分摊权重加载开销，（2）量化权重以减少每词元需要加载的字节数，以及（3）用推测解码在每步处理多个词元。

\translatornote{本题的 $6N$ 保留原书数值。常用近似是每个训练词元约需 $6N$ 次浮点运算，包含前向和反向传播。推理前向传播约为每词元 $2N$ 次，长上下文等情况还需计入注意力计算。参见 \href{https://arxiv.org/abs/2401.00448}{Sardana 等，2024，第 2 节}。}
\end{interview}
